\documentclass[10pt,a4paper]{amsart}
\usepackage{amsmath,amssymb,amsthm,graphicx}
% Retained minimal preamble for the Stacks Project chapter more-algebra.tex (More on Algebra).
% The chapter file pinned-original-source.tex carries only \input{preamble} on its first line, so no
% preamble text lives inside the chapter itself.  The Stacks Project uses one single shared preamble for
% every chapter; the snapshot stages it identically as sources/stacks-chow/preamble.tex and
% sources/stacks-etale-cohomology/preamble.tex.  The block below is copied mechanically, line for line,
% from sources/stacks-chow/preamble.tex lines 14-18 (the Xy-pic setup) and 154-217 (the native theorem
% declarations together with the mathematical macros).  Nothing was added, removed or redefined; the
% wrapper class and the amsmath/amssymb/amsthm/graphicx base packages are supplied by the assembler.
% Chapter: more-algebra.tex -- More on Algebra.
% For commutative diagrams we use Xy-pic
\usepackage[all]{xy}

% We use 2cell for 2-commutative diagrams.
\xyoption{2cell}

% Theorem environments.
%
\theoremstyle{plain}
\newtheorem{theorem}[subsection]{Theorem}
\newtheorem{proposition}[subsection]{Proposition}
\newtheorem{lemma}[subsection]{Lemma}

\theoremstyle{definition}
\newtheorem{definition}[subsection]{Definition}
\newtheorem{example}[subsection]{Example}
\newtheorem{exercise}[subsection]{Exercise}
\newtheorem{situation}[subsection]{Situation}

\theoremstyle{remark}
\newtheorem{remark}[subsection]{Remark}
\newtheorem{remarks}[subsection]{Remarks}

\numberwithin{equation}{subsection}

% Macros
%
\def\lim{\mathop{\mathrm{lim}}\nolimits}
\def\colim{\mathop{\mathrm{colim}}\nolimits}
\def\Spec{\mathop{\mathrm{Spec}}}
\def\Hom{\mathop{\mathrm{Hom}}\nolimits}
\def\Ext{\mathop{\mathrm{Ext}}\nolimits}
\def\SheafHom{\mathop{\mathcal{H}\!\mathit{om}}\nolimits}
\def\SheafExt{\mathop{\mathcal{E}\!\mathit{xt}}\nolimits}
\def\Sch{\mathit{Sch}}
\def\Mor{\mathop{\mathrm{Mor}}\nolimits}
\def\Ob{\mathop{\mathrm{Ob}}\nolimits}
\def\Sh{\mathop{\mathit{Sh}}\nolimits}
\def\NL{\mathop{N\!L}\nolimits}
\def\CH{\mathop{\mathrm{CH}}\nolimits}
\def\proetale{{pro\text{-}\acute{e}tale}}
\def\etale{{\acute{e}tale}}
\def\QCoh{\mathit{QCoh}}
\def\Ker{\mathop{\mathrm{Ker}}}
\def\Im{\mathop{\mathrm{Im}}}
\def\Coker{\mathop{\mathrm{Coker}}}
\def\Coim{\mathop{\mathrm{Coim}}}

% Boxtimes
%
\DeclareMathSymbol{\boxtimes}{\mathbin}{AMSa}{"02}

%
% Macros for moduli stacks/spaces
%
\def\QCohstack{\mathcal{QC}\!\mathit{oh}}
\def\Cohstack{\mathcal{C}\!\mathit{oh}}
\def\Spacesstack{\mathcal{S}\!\mathit{paces}}
\def\Quotfunctor{\mathrm{Quot}}
\def\Hilbfunctor{\mathrm{Hilb}}
\def\Curvesstack{\mathcal{C}\!\mathit{urves}}
\def\Polarizedstack{\mathcal{P}\!\mathit{olarized}}
\def\Complexesstack{\mathcal{C}\!\mathit{omplexes}}
% \Pic is the operator that assigns to X its picard group, usage \Pic(X)
% \Picardstack_{X/B} denotes the Picard stack of X over B
% \Picardfunctor_{X/B} denotes the Picard functor of X over B
\def\Pic{\mathop{\mathrm{Pic}}\nolimits}
\def\Picardstack{\mathcal{P}\!\mathit{ic}}
\def\Picardfunctor{\mathrm{Pic}}
\def\Deformationcategory{\mathcal{D}\!\mathit{ef}}

% --- appended 2026-10-12 (Chief): the project own definitions of the reference/slogan/history
% display environments, copied mechanically from the shared Stacks preamble lines 8-12
% (\usepackage{verbatim} plus the three \newenvironment lines, which the chapter preamble keeps
% verbatim from the project preamble). Nothing added, removed or redefined: these environments
% carry no mathematical content and the project itself typesets them as comments. Without these
% three lines a counted statement whose own text uses them fails with Environment slogan undefined.
\usepackage{verbatim}
\newenvironment{reference}{\comment}{\endcomment}
\newenvironment{slogan}{\comment}{\endcomment}
\newenvironment{history}{\comment}{\endcomment}
\usepackage{cleveref}
\crefname{originallemma}{original lemma}{original lemmas}
\Crefname{originallemma}{Original lemma}{Original lemmas}
\makeatletter
\newlabel{more-algebra-lemma-check-G-ring}{{07PQ}{source}{}{sourcepointer-1}{}}
\newlabel{more-algebra-lemma-check-G-ring@cref}{{[originallemma][0][]07PQ}{[1][0][]source}{}{}{}}
\makeatother
\crefname{originallemma}{original lemma}{original lemmas}
\Crefname{originallemma}{Original lemma}{Original lemmas}
\makeatletter
\newlabel{algebra-lemma-finite-free-given-residue-field-extension}{{07NE}{source}{}{sourcepointer-2}{}}
\newlabel{algebra-lemma-finite-free-given-residue-field-extension@cref}{{[originallemma][0][]07NE}{[1][0][]source}{}{}{}}
\makeatother
\crefname{originallemma}{original lemma}{original lemmas}
\Crefname{originallemma}{Original lemma}{Original lemmas}
\makeatletter
\newlabel{algebra-lemma-completion-finite-extension}{{07N9}{source}{}{sourcepointer-3}{}}
\newlabel{algebra-lemma-completion-finite-extension@cref}{{[originallemma][0][]07N9}{[1][0][]source}{}{}{}}
\makeatother
\crefname{originallemma}{original lemma}{original lemmas}
\Crefname{originallemma}{Original lemma}{Original lemmas}
\makeatletter
\newlabel{lemma-G-ring-goes-up-quasi-finite}{{07PP}{source}{}{sourcepointer-4}{}}
\newlabel{lemma-G-ring-goes-up-quasi-finite@cref}{{[originallemma][0][]07PP}{[1][0][]source}{}{}{}}
\makeatother
\title{A Noetherian ring is a G-ring exactly when every finite free extension has regular formal fibres}
\author{The Stacks Project Authors}
\date{stacks-project:more-algebra.tex:local-snapshot-20261012; GNU FDL 1.2 or later}
\begin{document}
\maketitle
\setcounter{subsection}{1}
\noindent\textit{Editorial note.} All delivered mathematical text is the author's own, copied from the pinned Stacks Project chapter file more-algebra.tex; no statement or proof step was rewritten, reordered or removed.\par\medskip
\noindent\textit{Editorial note.} The chapter file pinned-original-source.tex carries only \textbackslash{}input\{preamble\} on its first line, so no preamble text lives inside the chapter. The retained file sources/stacks-more-algebra/retained-preamble.tex is a mechanical, line-for-line copy of the project's single shared preamble as staged in the snapshot (Xy-pic setup plus theorem declarations and mathematical macros). Nothing was added, removed or redefined, and the excerpt serves only as preamble, never as a span coordinate.\par\medskip
\noindent\textit{Editorial note.} Both counted spans start at the \textbackslash{}begin line of their environment, so the statement and the author's proof are delivered whole; surrounding prose of the section is deliberately outside the counted spans and is carried as standard prerequisites instead.\par\medskip
\noindent\textit{Editorial note.} The citation fields (source\_id, source\_version, authors, source\_url, license.name, source\_title) are identical to those of every other card this lane has produced from this chapter, so the batch selector treats them as one source.\par\medskip
\noindent\textit{Allowed standard prerequisites (editorial).}\begin{itemize}
\item G-ring condition tested through formal fibres and purely inseparable residue field extensions. the discussion before the statement, source lines 12179-12261. Supplies the criterion the proof verifies.
\item Existence of a finite free extension realising a prescribed finite residue field extension at a prime lying over a given prime. the author cites Algebra, Lemma algebra-lemma-finite-free-given-residue-field-extension at proof lines 12275-12278. Produces the extension used in the computation.
\item Completion of a finite ring extension decomposes as a product of the completions at the primes lying over the given prime. the author cites Algebra, Lemma algebra-lemma-completion-finite-extension at proof lines 12279-12283. Gives the displayed product decomposition from which regularity is read off.
\end{itemize}\medskip
\ifdefined\hypertarget\hypertarget{sourcepointer-1}{}\fi
\noindent\textit{Source reference.} Original lemma 07PQ: focused-seeds/AGP/sources/stacks-snapshot-20261011/tags.txt, source lines 9805--9805. The lemma's own canonical tag, quoted so that readers can resolve the label inside the delivered text; it points at the counted statement itself and carries no proof content.\par\medskip
\ifdefined\hypertarget\hypertarget{sourcepointer-2}{}\fi
\noindent\textit{Source reference.} Original lemma 07NE: focused-seeds/AGP/sources/stacks-snapshot-20261011/tags.txt, source lines 9759--9759. Navigation pointer to a statement the counted proof refers to but does not contain; nothing load-bearing is delegated to it.\par\medskip
\ifdefined\hypertarget\hypertarget{sourcepointer-3}{}\fi
\noindent\textit{Source reference.} Original lemma 07N9: focused-seeds/AGP/sources/stacks-snapshot-20261011/tags.txt, source lines 9754--9754. Navigation pointer to a statement the counted proof refers to but does not contain; nothing load-bearing is delegated to it.\par\medskip
\ifdefined\hypertarget\hypertarget{sourcepointer-4}{}\fi
\noindent\textit{Source reference.} Original lemma 07PP: focused-seeds/AGP/sources/stacks-snapshot-20261011/tags.txt, source lines 9804--9804. Navigation pointer to a statement the counted proof refers to but does not contain; nothing load-bearing is delegated to it.\par\medskip
\begin{lemma}
\label{lemma-check-G-ring}
Let $R$ be a Noetherian ring. Then $R$ is a G-ring if and only if
for every finite free ring map $R \to S$ the formal fibres of $S$
are regular rings.
\end{lemma}

\begin{proof}
Assume that for any finite free ring map $R \to S$ the ring $S$ has
regular formal fibres. Let $\mathfrak q \subset \mathfrak p \subset R$
be primes and let $\kappa(\mathfrak q) \subset L$ be a finite purely
inseparable extension. To show that $R$ is a G-ring it suffices to
show that
$$
R_\mathfrak p^\wedge \otimes_R \kappa(\mathfrak q)
\otimes_{\kappa(\mathfrak q)} L
$$
is a regular ring. Choose a finite free extension $R \to R'$ such that
$\mathfrak q' = \mathfrak qR'$ is a prime and such that $\kappa(\mathfrak q')$
is isomorphic to $L$ over $\kappa(\mathfrak q)$, see
Algebra, Lemma \ref{algebra-lemma-finite-free-given-residue-field-extension}.
By
Algebra, Lemma \ref{algebra-lemma-completion-finite-extension}
we have
$$
R_\mathfrak p^\wedge \otimes_R R' = \prod (R'_{\mathfrak p_i'})^\wedge
$$
where $\mathfrak p_i'$ are the primes of $R'$ lying over $\mathfrak p$.
Thus we have
$$
R_\mathfrak p^\wedge \otimes_R \kappa(\mathfrak q)
\otimes_{\kappa(\mathfrak q)} L =
R_\mathfrak p^\wedge \otimes_R R'
\otimes_{R'} \kappa(\mathfrak q')
=
\prod (R'_{\mathfrak p_i'})^\wedge
\otimes_{R'_{\mathfrak p'_i}} \kappa(\mathfrak q')
$$
Our assumption is that the rings on the right are regular, hence the
ring on the left is regular too. Thus $R$ is a G-ring. The converse
follows from Lemma \ref{lemma-G-ring-goes-up-quasi-finite}.
\end{proof}

\end{document}
